\begin{table}[H]
\centering
\caption{Remuneración de ocupados de 25 a 29 años con pregrado universitario, 2021 y 2025}
\label{tab:remuneracion_pregrado_universitario_25_29}
\footnotesize
\begin{tabular}{@{}p{7.4cm}rr@{}}
\toprule
Indicador & 2021 & 2025 \\
\midrule
\multicolumn{3}{@{}l}{\textbf{Panel A. Remuneración mensual por trabajador}} \\
Ocupados (miles) & 347,1 & 454,7 \\
Promedio (millones de pesos de 2025) & 3,1 & 2,8 \\
Mediana (millones de pesos de 2025) & 2,5 & 2,5 \\
Desviación estándar (millones de pesos de 2025) & 3,6 & 1,9 \\
Razón P90/P10 & 4,9 & 3,6 \\
Entre 0,9 y 1,1 SMLMV & 8,8\% & 13,3\% \\
\addlinespace[0.6em]
\multicolumn{3}{@{}l}{\textbf{Panel B. Remuneración por hora trabajada}} \\
Promedio (miles de pesos de 2025) & 16,4 & 15,6 \\
Mediana (miles de pesos de 2025) & 11,8 & 13,0 \\
Desviación estándar (miles de pesos de 2025) & 22,1 & 11,5 \\
Razón P90/P10 & 4,8 & 3,9 \\
\bottomrule
\end{tabular}
\caption*{\footnotesize Nota: la muestra se restringe a ocupados de 25 a 29 años con \texttt{P3042} igual a universitaria y \texttt{P3043} igual a título universitario. El SMLMV se expresa en pesos constantes de 2025 y no incluye auxilio de transporte. Fuente: cálculos propios con microdatos mensuales de la GEIH marco 2018 del DANE.}
\end{table}
